\label{sec:introduction}

Large language models (LLMs) increasingly answer questions by combining
parametric knowledge with retrieved documents, tools, memory and long
conversational histories. These mechanisms improve evidence coverage, but they
do not guarantee that a system identifies the \emph{situation} in which a fact
is valid. A model may read that a person was appointed and later resigned, yet
still name that person as the current office holder. It may turn a proposal
into a decision, extend a local policy to an entire organization, answer from
information unavailable to the queried observer, or collapse a hypothetical
scenario into the real world. In each case the decisive evidence can already be
present. The failure is not simply missing knowledge; it is incorrect evidence
applicability.

We define \emph{situation blindness} as producing an answer or action that is
supported by at least one available claim but inconsistent with the time, scope,
modality, observer knowledge or possible world required by the query. This
category matters because conventional factuality checks can mark the supporting
claim as true while missing that it is true in the wrong situation. Natural
question-answering resources already show that a material subset of ordinary
questions depends on extra-linguistic time or place, and that updated evidence
does not eliminate the resulting performance loss~\cite{zhang2021situatedqa,
vu2023freshqa,meem2024pat,chen2021timeqa,chen2025chronoqa}.

Retrieval-augmented generation supplies external evidence~\cite{lewis2020rag};
self-reflective and tool-using systems add critique and action
~\cite{asai2024selfrag,yao2023react,schick2023toolformer}; context engineering
organizes the inference-time information payload~\cite{mei2025context}; and
agent harnesses mediate tools, memory, permissions, verification and feedback
~\cite{zhong2026harness,lin2026ahe}. These layers are necessary, but none
requires an explicit answer to a separate question: \emph{which interpretation
of the available evidence governs this interaction now?}

We call the corresponding discipline \emph{situation engineering}: the
systematic design of representations, sensing interfaces, update rules,
validation procedures and response policies that establish the active
query-relative world state before an AI system commits to an answer or action.
For a query $q$ and evidence $E$, a minimal state is
\begin{equation}
S(q,E)=\{\mathcal{A},\mathcal{V},\mathcal{T},\mathcal{P},\mathcal{M},
\mathcal{K},\mathcal{W},\mathcal{U}\},
\end{equation}
where the variables denote actors, events, valid time, scope, modality and
source status, observer-specific knowledge, possible world and unresolved
answer-critical information. A policy then selects
\begin{equation}
\pi(S)\in\{\textsc{Answer},\textsc{Clarify},\textsc{Retrieve},
\textsc{Abstain},\textsc{Act}\}.
\end{equation}
This decomposition separates evidence availability from evidence applicability.

We organize the study around three questions.
\textbf{RQ1 (representation):} can failures involving time, scope, modality,
source status, observer knowledge, possible world and missing information be
expressed through one inspectable applicability interface?
\textbf{RQ2 (localization):} when that interface fails, can matched controls
separate state-sensing errors from update, validation and response-policy
errors?
\textbf{RQ3 (transfer):} does exposing the interface to current LLMs improve
answers beyond equally informed structured, retrieval and tool-using
baselines? RQ1 and RQ2 are evaluated as controlled software and mechanistic
claims. RQ3 is a deliberately bounded pilot and remains open as a general
empirical claim.

This question structure yields four contributions with different evidential
status. First, we operationalize \emph{situation blindness}: correctness
requires both evidential support and query-relative applicability. Second, we
specify a systems contract connecting sensing, provenance-linked state,
validation and answer/retrieve/clarify/abstain/action policy, while leaving the
choice of logical, probabilistic or neural implementation open. Third, we
release SCQA and SituationCatch-Bench as mechanism-level diagnostics: matched
state conditions and component ablations test whether each declared variable
has its predicted behavioural consequence. Fourth, we contribute a negative
and therefore informative boundary result. In an audited three-model
Gemini-family pilot, an explicit situation prompt is null-to-negative against a
matched structured prompt and trails focused retrieval overall. This prevents
the framework contribution from being conflated with prompt superiority and
pre-specifies the natural-text, multi-family and human-agreement evidence
needed to test transfer.
